% Preámbulo de la versión en español.

\usepackage[utf8]{inputenc}
\usepackage[T1]{fontenc}
\usepackage[spanish,es-noshorthands,es-tabla]{babel}

\usepackage[a4paper,top=2.5cm,bottom=2.5cm,left=2.5cm,right=2.5cm]{geometry}
\usepackage{setspace}
\onehalfspacing

\usepackage{amsmath,amssymb,amsthm,mathtools}
\usepackage{bm}

\usepackage{graphicx}
\usepackage{xcolor}
\usepackage{tikz}
\usetikzlibrary{arrows.meta,positioning,calc}
\usepackage{enumitem}
\usepackage{booktabs}
\usepackage{caption}
\usepackage[nottoc]{tocbibind}
\usepackage[colorlinks=true,linkcolor=black,citecolor=blue,urlcolor=blue]{hyperref}
\hypersetup{
  pdftitle={Mecánica cuántica euclídea aplicada a la valoración de opciones path-dependent (borrador)},
  pdfauthor={Roberto Cepeda Rocandio}
}

\theoremstyle{plain}
\newtheorem{teorema}{Teorema}[chapter]
\newtheorem{proposicion}[teorema]{Proposición}
\newtheorem{lema}[teorema]{Lema}
\newtheorem{corolario}[teorema]{Corolario}

\theoremstyle{definition}
\newtheorem{definicion}[teorema]{Definición}
\newtheorem{ejemplo}[teorema]{Ejemplo}

\theoremstyle{remark}
\newtheorem{observacion}[teorema]{Observación}

\newcommand{\R}{\mathbb{R}}
\newcommand{\N}{\mathbb{N}}
\newcommand{\Z}{\mathbb{Z}}
\newcommand{\C}{\mathbb{C}}
\newcommand{\E}{\mathbb{E}}
\newcommand{\PP}{\mathbb{P}}
\newcommand{\Q}{\mathbb{Q}}
\newcommand{\F}{\mathcal{F}}
\newcommand{\G}{\mathcal{G}}
\newcommand{\Hh}{\mathcal{H}}
\newcommand{\HH}{\mathcal{H}}
\newcommand{\1}{\mathbf{1}}
\newcommand{\dd}{\mathrm{d}}
\newcommand{\eqdef}{\coloneqq}
\DeclareMathOperator{\Var}{Var}

\newcommand{\bra}[1]{\langle #1 |}
\newcommand{\ket}[1]{| #1 \rangle}
\newcommand{\braket}[2]{\langle #1 \mid #2 \rangle}
\newcommand{\matel}[3]{\langle #1 \mid #2 \mid #3 \rangle}
\newcommand{\Hbs}{H_{\text{BS}}}
\newcommand{\Heff}{H_{\text{eff}}}
\newcommand{\Lag}{\mathcal{L}}
\newcommand{\Dpath}{\mathcal{D}}

% Borrador en curso: remisiones y capítulos aún no publicados.
\newcommand{\enprep}{, en preparación}
\newif\ifprimerprevisto \primerprevistotrue
\newcommand{\capituloprevisto}[2]{%
  \ifprimerprevisto\clearpage\global\primerprevistofalse\else\par\vspace{2.5em}\fi
  \begingroup\let\clearpage\relax\let\cleardoublepage\relax
  \chapter{#1}\label{#2}\endgroup
  \noindent\textit{Capítulo en preparación.}\par}
\newcommand{\seccionesprevistas}[1]{%
  \par\medskip\noindent\textit{Secciones #1: en preparación.}\par}

\usepackage{titlesec}
\titleformat{\chapter}[hang]{\normalfont\huge\bfseries}{\thechapter.}{0.6em}{}
\titlespacing*{\chapter}{0pt}{-20pt}{18pt}
